\chapter{Divisor Substitutions and Satake Parameters}
\label{v4-arith:kazhdan-beilinson}

The divisor substitutions on polynomial Fock space, Adams ring maps,
and a spherical Hecke representation are different operator data.
We specify the Satake normalization, compute the divisor discrepancy,
and distinguish a Satake matrix from its scalar symmetric-power traces.

\section{The Satake Pair and Why It Matters}
\label{v4-arith:kb:satake-pair}

\begin{theorem}[normalized spherical Satake convention]
\label{v4-arith:thm:satake-GL2}
Let \(K=GL_2(\mathbb Z_p)\), normalize Haar measure by
\(\operatorname{vol}(K)=1\), and use complex coefficients.
The normalized spherical Satake isomorphism identifies
\(\mathcal H_p=C_c(K\backslash GL_2(\mathbb Q_p)/K,\mathbb C)\)
with \(\mathbb C[X^{\pm1},Y^{\pm1}]^{S_2}\).
The characteristic functions of \(K\operatorname{diag}(p,1)K\)
and \(K\operatorname{diag}(p,p)K\), denoted \(T_p\) and \(D_p\),
have images \(p^{1/2}(X+Y)\) and \(XY\).
For a character of this algebra set
\(\alpha_p=p^{1/2}X\), \(\beta_p=p^{1/2}Y\) at that character.
Then \(\alpha_p\beta_p=pD_p\), where the right side means its
scalar value. Trivial central character gives \(D_p=1\).
The condition \(|X|=|Y|=1\) is the additional tempered condition,
not a consequence of trivial central character alone. These conventions
are Gross, \emph{On the Satake isomorphism}, Proposition3.6,
equations(3.13)--(3.14) and \S8.
\end{theorem}

\begin{proposition}[symmetric-power generating function]
\label{v4-arith:prop:hecke-generating-function}
Define the elements \(H_{p^k}\in\mathcal H_p\) by their Satake images
\(h_k(p^{1/2}X,p^{1/2}Y)\), the complete symmetric polynomials.
Then \(H_{p^0}=1\), \(H_p=T_p\), and
\[
 \sum_{k\ge0}H_{p^k}t^k=(1-T_pt+pD_pt^2)^{-1}.
\]
\end{proposition}
\begin{proof}
In the Satake algebra the generating series is
\((1-p^{1/2}Xt)^{-1}(1-p^{1/2}Yt)^{-1}\).
Multiply the denominators and use injectivity of the Satake map.
This fixes the determinant normalization and the choice of the
higher Hecke elements rather than identifying them with a single
undeclared double-coset convention.
\end{proof}

At a character with \(D_p=1\) this series is
\(\bigl((1-\alpha_pt)(1-\beta_pt)\bigr)^{-1}\).
The two parameters may coincide. In the tempered case their unitary
versions \(p^{-1/2}\alpha_p,p^{-1/2}\beta_p\) give the local
trigonometric Euler factor; changing between the two normalizations
shifts the Dirichlet variable by one half.
This complex Satake convention is separate from the integral polynomial
carrier used for the divisor substitutions below.

\section{Adams operations and Hecke representations}
\label{v4-arith:kb:adams-not-hecke}

\begin{proposition}
\label{v4-arith:prop:adams-commutative-monoid}
The Adams operations of a special $\lambda$-ring
(Definition~\ref{v4-arith:hecke:def:lambda}) satisfy
$\psi_m\psi_n=\psi_{mn}$.
They give an action of $(\mathbb N_{\geq1},\times)$,
the free commutative monoid on the primes.
Conversely, commuting Frobenius lifts determine the special structure
under the torsion-free hypothesis of
Theorem~\ref{v4-arith:hecke:thm:wilkerson}.
Example~\ref{v4-arith:hecke:ex:torsion-lambda} shows why that
hypothesis cannot be omitted.
This action need not extend to the group completion: the polynomial
index shifts are not surjective for \(n>1\).

The full local spherical \(GL_1\) Hecke algebra over \(\mathbb C\)
is \(\mathbb C[\mathbb Z]=\mathbb C[z,z^{-1}]\).
Valuation identifies the double cosets of
\(\mathbb Q_p^\times/\mathbb Z_p^\times\) with \(\mathbb Z\), and,
with the unit subgroup of volume one, convolution adds valuations.
An Adams action supplies only the multiplicative monoid action;
a comparison with a specified positive Hecke subalgebra is separate
data.
\end{proposition}

\begin{observation}[distinct operator data]
\label{v4-arith:obs:gl1-vs-gl2}
An Adams action is an action of the multiplicative monoid by ring
endomorphisms. It does not by itself specify a \(GL_2\) Hecke
representation, its Satake parameters, or its determinant operation.
A \(\Lambda\)-ring can contain symmetric functions in two parameters;
the missing datum is their comparison with the specified arithmetic
operators, not an impossibility of encoding an unordered pair.
\end{observation}

\section{Explicit Computation: The Recursion Fails}
\label{v4-arith:kb:recursion-fails}

Let \(\mathcal{H}^{(p)}_{\mathrm{rk1}}\) denote the rank-\(1\) Heisenberg
algebra at prime \(p\), with mode operators \(a^{(p)}_{n}\) (\(n\in\mathbb{Z}\))
satisfying \([a^{(p)}_{m},a^{(p)}_{n}]=m\,\delta_{m+n,0}\), acting
on the Fock space \(\mathcal{F}^{(p)}\) generated from the vacuum
\(|0\rangle\) (annihilated by \(a^{(p)}_{n}\) for \(n\geq 0\)).
For each \(k\ge1\), the state \(a^{(p)}_{-k}|0\rangle\) spans the
one-particle subspace of weight \(k\). The full weight-\(k\) subspace
also contains the multi-particle partition states.

The claimed formula for the action of a putative \(GL_{2}\) Hecke
operator \(T_{n}\) on this Fock space is
\begin{equation}
T_{n}(a^{(p)}_{-k}|0\rangle)\;=\;\sum_{d\mid n}d\cdot a^{(p)}_{-nk/d}|0\rangle\cdot\mathbf{1}_{d\mid k},
\end{equation}
extended by \(T_{n}(|0\rangle)=|0\rangle\).

\begin{computation}
\label{v4-arith:comp:T_p-on-a_-1}
\(T_{p}(a^{(p)}_{-1}|0\rangle)\). Divisors of \(n=p\) are
\(d\in\{1,p\}\). Only \(d=1\) satisfies \(d\mid k=1\). So
\begin{equation}
T_{p}(a^{(p)}_{-1}|0\rangle)\;=\;a^{(p)}_{-p}|0\rangle.
\end{equation}
\end{computation}

\begin{computation}
\label{v4-arith:comp:T_p-squared}
\(T_{p}^{2}(a^{(p)}_{-1}|0\rangle)\). Apply \(T_{p}\) to
\(a^{(p)}_{-p}|0\rangle\): divisors \(d\in\{1,p\}\), both satisfy
\(d\mid p\), so
\begin{equation}
T_{p}(a^{(p)}_{-p}|0\rangle)\;=\;a^{(p)}_{-p^{2}}|0\rangle
+p\cdot a^{(p)}_{-p}|0\rangle.
\end{equation}
Therefore
\begin{equation}
T_{p}^{2}(a^{(p)}_{-1}|0\rangle)\;=\;a^{(p)}_{-p^{2}}|0\rangle
+p\cdot a^{(p)}_{-p}|0\rangle.
\end{equation}
\end{computation}

\begin{computation}
\label{v4-arith:comp:T_p-squared-direct}
\(T_{p^{2}}(a^{(p)}_{-1}|0\rangle)\). Divisors \(d\in\{1,p,p^{2}\}\),
only \(d=1\) satisfies \(d\mid 1\). So
\begin{equation}
T_{p^{2}}(a^{(p)}_{-1}|0\rangle)\;=\;a^{(p)}_{-p^{2}}|0\rangle.
\end{equation}
\end{computation}

\begin{observation}[recursion discrepancy]
\label{v4-arith:obs:recursion-discrepancy}
Testing \(T_{p^{2}}=T_{p}^{2}-p\cdot\mathrm{id}\) on \(a^{(p)}_{-1}|0\rangle\):
\begin{align}
T_{p}^{2}-p\cdot\mathrm{id}\bigr|_{a^{(p)}_{-1}|0\rangle}
&\;=\;a^{(p)}_{-p^{2}}|0\rangle+p\cdot a^{(p)}_{-p}|0\rangle-p\cdot a^{(p)}_{-1}|0\rangle, \\
T_{p^{2}}\bigr|_{a^{(p)}_{-1}|0\rangle}
&\;=\;a^{(p)}_{-p^{2}}|0\rangle.
\end{align}
The discrepancy is \(p(a^{(p)}_{-p}-a^{(p)}_{-1})|0\rangle\neq 0\).
\end{observation}

\begin{corollary}[falsification]
\label{v4-arith:cor:recursion-falsified}
The formula
\(T_{n}(a^{(p)}_{-k})=\sum_{d\mid n}d\cdot a^{(p)}_{-nk/d}\mathbf{1}_{d\mid k}\)
fails both universal Adams multiplicativity and the displayed
\(GL_2\) Hecke recursion. It satisfies coprime multiplicativity and
the prime-valuation composition law of
Corollary~\ref{v4-arith:ado:global-composition}. The preceding computations
prove the two failures already on the first generator.
\end{corollary}

\section{Diagonalisation onto the Satake Pair}
\label{v4-arith:kb:correction-satake-diagonalisation}

A two-dimensional Satake model gives a different construction.
Choose nonzero complex numbers \(\alpha_p,\beta_p\) with
\(\alpha_p\beta_p=p\). This is parameter data, not a consequence of the
rank of an oscillator algebra. In the tempered parametrization one can
write \(\alpha_p=p^{1/2}e^{i\theta_p}\),
\(\beta_p=p^{1/2}e^{-i\theta_p}\) with real \(\theta_p\).

\begin{definition}[Satake operator]
\label{v4-arith:def:satake-diagonalised}
Let \(F_p^1=\mathbb C a^{(p)}_{-1}|0\rangle\oplus
\mathbb C b^{(p)}_{-1}|0\rangle\). Define
\(S_p=\operatorname{diag}(\alpha_p,\beta_p)\) in this basis.
The equation \(\alpha_p\beta_p=p\) is imposed on the parameters.
It is not encoded by multiplying a Heisenberg conformal vector by
\(\log p\); logarithmic energy is the separate convention of
Convention~\ref{v4-arith:conv:two-grading}.
\end{definition}

\begin{theorem}[trace of the Satake operator]
\label{v4-arith:thm:corrected-hecke}
The operator \(S_p\) has determinant \(p\) and trace
\(\alpha_p+\beta_p\), equal to \(2p^{1/2}\cos\theta_p\) in the tempered
parametrization. The line spanned by
\((a^{(p)}_{-1}+b^{(p)}_{-1})|0\rangle\) is invariant precisely when
\(\alpha_p=\beta_p\).
\end{theorem}
\begin{proof}
The determinant and trace are those of the displayed diagonal matrix.
Its value on the sum of the two basis vectors is proportional to that
sum exactly when its two diagonal entries agree.
\end{proof}

\begin{corollary}[symmetric-power trace recursion]
\label{v4-arith:cor:hecke-recursion-Sym}
For \(k\ge0\), set \(h_k=\operatorname{Tr}(\operatorname{Sym}^kS_p)\).
Then \(h_0=1\), \(h_1=\alpha_p+\beta_p\), and
\[
 h_k=\sum_{j=0}^k\alpha_p^{k-j}\beta_p^j,\qquad
 h_{k+1}=(\alpha_p+\beta_p)h_k-p h_{k-1}\quad(k\ge1).
\]
Thus the scalar operators \(H_{p^k}=h_k\,\mathrm{id}\) on a fixed
one-dimensional space satisfy the displayed local Hecke recursion.
This is a scalar-parameter representation, distinct from \(S_p\) on
\(F_p^1\) and from the divisor substitutions \(T_{p^k}\).
\end{corollary}
\begin{proof}
The monomial basis of \(\operatorname{Sym}^kF_p^1\) has eigenvalues
\(\alpha_p^{k-j}\beta_p^j\). Summing gives \(h_k\). Multiplying this
sum by \(\alpha_p+\beta_p\) and subtracting
\(\alpha_p\beta_p h_{k-1}\) leaves exactly the sum for \(h_{k+1}\).
\end{proof}

\section{The Angle \(\theta_{p}\) as Unramified Langlands Parameter}
\label{v4-arith:kb:theta-p-background}

The angle \(\theta_{p}\) first appeared in Definition~\ref{v4-arith:def:satake-diagonalised}
as the argument of the Satake parameters \(\alpha_{p}=p^{1/2}e^{i\theta_{p}}\),
\(\beta_{p}=p^{1/2}e^{-i\theta_{p}}\). The question is whether \(\theta_{p}\) is
a free internal parameter of the Heisenberg algebra or is fixed by external
(Langlands) data. For a specified unramified automorphic representation it is determined
by its Satake pair, when that pair has the tempered normalization above.
An arbitrary choice of angles does not construct such a representation.

\begin{proposition}[tempered determinant normalization]
\label{v4-arith:prop:theta-p-is-satake}
If \(\alpha_p\beta_p=p\) and
\(|\alpha_p|=|\beta_p|=p^{1/2}\), then the unordered pair is
\(\{p^{1/2}e^{i\theta_p},p^{1/2}e^{-i\theta_p}\}\), where
\(\theta_p\) is defined modulo \(2\pi\) and sign.
\end{proposition}
\begin{proof}
The number \(\alpha_p/p^{1/2}\) has modulus one, so it is
\(e^{i\theta_p}\). The determinant condition forces the second
normalized number to be \(e^{-i\theta_p}\). Interchanging the two
entries changes the sign of \(\theta_p\).
\end{proof}

\begin{convention}
\(\{\theta_{p}\}_{p}\) is chosen once and for all as a map
\(\theta\colon\mathbb P\to\mathbb R/(2\pi\mathbb Z,\,\theta\sim-\theta)\)
(the angle modulo \(\theta\leftrightarrow-\theta\), since
\(2\cos\theta=2\cos(-\theta)\)). This is the data of an
unramified local Langlands parameter at each finite place. The
Hartnoll--Yang L-function
\begin{equation}
L(s)\;=\;\prod_{p}\frac{1}{1-2\cos\theta_{p}\,p^{-s}+p^{-2s}}
\end{equation}
is a formal Euler product. It is the standard \(L\)-function of an
automorphic representation only if that representation and the comparison
with its normalized local parameters are supplied; arbitrary angles do
not establish automorphy or analytic continuation.
\end{convention}

\section{A radial Hecke comparison}
\label{v4-arith:radial:section}

The Satake trace recursion also admits a representation on radial functions
on a regular tree.
Its adjacency operator has the density used in
Definition~\ref{v4-arith:def:averaging-naive}. Matching that density
does not identify the operators. The radial tree representation gives an
explicit comparison with the spherical Hecke action and its vector state.
This is the classical spherical transform, with the distance and
symmetric-power normalizations kept separate. The density is the
Kesten--McKay measure. McKay, \emph{The expected eigenvalue distribution
of a large regular graph}, Linear Algebra Appl. 40 (1981),
203--216, Theorem~1.1, gives its graph-theoretic normalization.

Fix a prime $p$ and a rooted $(p+1)$-regular tree $(\mathcal T_p,o)$.
Use counting measure on its vertices and complex coefficients.
Let $S_n$ be the vertices at distance $n$ from $o$, and write
\[
 s_0=1,\qquad s_n=(p+1)p^{n-1}\quad(n\ge1),\qquad
 e_n=s_n^{-1/2}\mathbf{1}_{S_n}.
\]
The radial Hilbert space $\mathcal R_p$ consists of the square-summable
functions constant on each $S_n$. The vectors $e_n$ form an orthonormal
basis. All these vector spaces have cohomological degree zero.
Distance gives a filtration, not a grading preserved by adjacency.

Let $T$ be adjacency, $(Tf)(v)=\sum_{w\sim v}f(w)$.
It is a bounded self-adjoint operator on the vertex Hilbert space,
with norm at most $p+1$ by Cauchy--Schwarz.
Counting edges between consecutive spheres gives
\begin{align}
 Te_0&=\sqrt{p+1}\,e_1,\\
 Te_1&=\sqrt{p+1}\,e_0+\sqrt p\,e_2,\\
 Te_n&=\sqrt p\,(e_{n-1}+e_{n+1})\quad(n\ge2).
 \label{v4-arith:radial:recurrence}
\end{align}
In particular, two steps return to the root in $p+1$ ways.
Thus the polynomial which produces the distance-two sphere is
$T^2-(p+1)I$. The symmetric-power Hecke polynomial is instead
$T^2-pI$. Their difference is the identity operator.

Define monic integral polynomials by
\[
 R_0(x)=1,\quad R_1(x)=x,\quad R_2(x)=x^2-(p+1),\qquad
 R_{n+1}(x)=xR_n(x)-pR_{n-1}(x)\quad(n\ge2).
\]
The preceding recurrence proves
$R_n(T)e_0=\mathbf{1}_{S_n}$. These formulas already give an integral
$\mathbb Z[x]$-module isomorphism
\[
 \mathbb Z[x]\longrightarrow
 \bigoplus_{n\ge0}\mathbb Z\mathbf{1}_{S_n},\qquad
 f\longmapsto f(T)e_0.
\]
Here $x$ acts on the right by $T$. Monicity makes the $R_n$ an integral
basis, so this assertion needs no inversion of $p$.
The normalized Hilbert comparison uses real square roots and complex scalars.

\begin{theorem}[radial spherical transform]
\label{v4-arith:radial:transform}
Put $x(\theta)=2\sqrt p\cos\theta$ for $0\le\theta\le\pi$, and set
\[
 d\nu_p(\theta)=\frac{2p(p+1)\sin^2\theta}
 {\pi(p^2+1-2p\cos2\theta)}\,d\theta,\qquad
 q_n(\theta)=\frac{R_n(x(\theta))}{\sqrt{s_n}}.
\]
The assignment $\mathcal U_pe_n=q_n$ extends uniquely to a unitary
\[
 \mathcal U_p:\mathcal R_p\longrightarrow L^2([0,\pi],\nu_p).
\]
It sends $e_0$ to $1$ and satisfies the commutative square
\[
 \begin{array}{ccc}
 \mathcal R_p &\xrightarrow{T}& \mathcal R_p\\
 \mathcal U_p\downarrow &&\downarrow\mathcal U_p\\
 L^2(\nu_p)&\xrightarrow{M_x}&L^2(\nu_p)
 \end{array}
\]
where $M_x$ is multiplication by $x(\theta)$.
Consequently $\|T|_{\mathcal R_p}\|=2\sqrt p$, and for every bounded
Borel function $f$ on $[-2\sqrt p,2\sqrt p]$,
\[
 \langle e_0,f(T)e_0\rangle
 =\int_0^\pi f(2\sqrt p\cos\theta)\,d\nu_p(\theta).
 \label{v4-arith:radial:state}
\]
\end{theorem}

\begin{proof}
The tridiagonal recurrence makes $e_0$ cyclic for $T$ on $\mathcal R_p$.
We compute its spectral measure to determine the Hilbert norm of the
polynomial map. Remove the first row and column of that tridiagonal
matrix. Its tail has every off-diagonal coefficient equal to $\sqrt p$.
For $\operatorname{Im}z>0$, the first diagonal entry $g(z)$ of its
resolvent $(z-J)^{-1}$ satisfies the Schur-complement identity
\[
 g(z)=\frac1{z-pg(z)}
      =\frac{z-\sqrt{z^2-4p}}{2p}.
\]
The square root is analytic off $[-2\sqrt p,2\sqrt p]$ and is asymptotic
to $z$ at infinity. The asymptotic $g(z)\sim z^{-1}$ selects this root.
The root resolvent therefore equals
\[
 m(z)=\langle e_0,(z-T)^{-1}e_0\rangle
      =\frac1{z-(p+1)g(z)}.
\]
For $|x|<2\sqrt p$, substitution and Stieltjes inversion give
\[
 -\frac1\pi\operatorname{Im}m(x+i0)
 =\frac{(p+1)\sqrt{4p-x^2}}
 {2\pi((p+1)^2-x^2)}.
\]
There are no poles off this interval. A zero of the denominator would
force $g^2=1$ and $z=\pm(p+1)$, whereas the chosen branch has
$g(\pm(p+1))=\pm p^{-1}$. At both endpoints the denominator is
nonzero because $p>1$. Thus $m$ has continuous finite boundary values
on the interval, including its endpoints, and analytic continuation
across its complement. The inversion formula gives only the displayed
absolutely continuous measure. Its total mass is one, since
$m(z)\sim z^{-1}$. Changing variables $x=2\sqrt p\cos\theta$
gives $\nu_p$.

The spectral theorem now gives
\[
 \int R_n(x)R_m(x)\,d\nu_p
 =\langle R_n(T)e_0,R_m(T)e_0\rangle
 =s_n\delta_{nm}.
\]
Hence the $q_n$ are orthonormal. Polynomials are dense in continuous
functions on the compact $x$ interval and thus in its $L^2$ space.
The monic $R_n$ span all polynomials, so the $q_n$ are complete.
This proves that $\mathcal U_p$ is unitary. The recurrence proves
intertwining first on finite sums and then everywhere by boundedness.
The support of the measure is the full closed interval. Multiplication
by $x$ therefore has norm $2\sqrt p$. Bounded Borel functional calculus
proves the state identity.
\end{proof}

\section{Hecke operations and the distinguished state}
\label{v4-arith:radial:hecke}

For the Bruhat--Tits tree of $\mathrm{PGL}_2(\mathbb Q_p)$, vertices are
homothety classes of $\mathbb Z_p$-lattices in $\mathbb Q_p^2$.
Two vertices are adjacent when representatives satisfy
$pL\subset L'\subset L$ with $[L:L']=p$.
The lines in $L/pL$ give exactly $p+1$ neighbors.
With the maximal compact subgroup of Haar volume one, its adjacency
operator is the characteristic-function Hecke operator for
$\operatorname{diag}(p,1)$. Pullback to $\mathrm{GL}_2$ fixes the
central operator $\operatorname{diag}(p,p)$ to be the identity.
Thus the representation has trivial central character.

Define $H_0(x)=1$, $H_1(x)=x$ and
$H_{k+1}(x)=xH_k(x)-pH_{k-1}(x)$ for $k\ge1$.
In the notation of Proposition~\ref{v4-arith:prop:hecke-generating-function},
$H_k(T)$ represents $H_{p^k}$ with $D_p=1$. These are the
complete-symmetric Satake elements in this central character. Their generating series and spectral values are
\[
 \sum_{k\ge0}H_k(x)t^k=(1-xt+pt^2)^{-1},\qquad
 H_k(x(\theta))=p^{k/2}\sum_{j=0}^ke^{i(k-2j)\theta}.
\]
The second formula follows by multiplying the two geometric series
with parameters $\sqrt p\,e^{i\theta}$ and $\sqrt p\,e^{-i\theta}$.
It remains valid at $\theta=0,\pi$, where the parameters coincide.
The transform therefore intertwines every $H_k(T)$ with the stated
multiplication operator, with the same cyclic state.

\begin{proposition}[distance and symmetric-power elements]
\label{v4-arith:radial:distance}
Let $(D_nf)(v)=\sum_{d(v,w)=n}f(w)$ on the vertex Hilbert space.
Then $D_n=R_n(T)$ and
\[
 H_k(T)=\sum_{j=0}^{\lfloor k/2\rfloor}D_{k-2j},\qquad
 \langle e_0,H_k(T)e_0\rangle=
 \begin{cases}1&k\text{ even},\\0&k\text{ odd}.\end{cases}
\]
On $\mathcal R_p$, for $|t|<p^{-1/2}$, the series converges in operator norm and
\[
 \left\langle e_0,(I-tT+pt^2I)^{-1}e_0\right\rangle
 =\frac1{1-t^2}.
\]
\end{proposition}
\begin{proof}
The distance operators are bounded because every row and column has
$s_n$ entries. Counting the first step gives
$TD_1=D_2+(p+1)I$ and
$TD_n=D_{n+1}+pD_{n-1}$ for $n\ge2$. Thus $D_n=R_n(T)$.
With $H_{-1}=H_{-2}=0$, the recurrences give
$R_n=H_n-H_{n-2}$ for every $n\ge0$.
Telescoping proves the sum formula. Since $D_ne_0=\mathbf{1}_{S_n}$,
only $D_0$ contributes to the root state.
The spectral formula bounds $\|H_k(T)|_{\mathcal R_p}\|$
by $(k+1)p^{k/2}$. This proves norm convergence on $\mathcal R_p$
in the stated disk and permits evaluation of its state term by term.
Multiplication of the series by $I-tT+pt^2I$ proves the inverse formula
on that space.
\end{proof}

For example, at $p=2$,
\[
 q_0=1,\qquad q_1=x/\sqrt3,\qquad
 q_2=(x^2-3)/\sqrt6,\qquad H_2=x^2-2.
\]
The tree counts give
$\int x^2d\nu_2=3$ and $\int x^4d\nu_2=15$.
Indeed a four-step return either returns after two steps, giving nine
walks, or first reaches distance two, giving six walks.
Consequently $\int H_2d\nu_2=1$ but $\int R_2d\nu_2=0$,
and $\int R_2^2d\nu_2=6$. These values distinguish the two Hecke
normalizations already at the first nontrivial higher element.

\section{Finite sets of primes and the arithmetic comparison}
\label{v4-arith:radial:finite}

For a finite prime set $F$, form the completed Hilbert tensor products
$\mathcal R_F=\widehat\bigotimes_{p\in F}\mathcal R_p$ and the product
probability measure $\nu_F=\prod_{p\in F}\nu_p$.
The tensor product of the maps $\mathcal U_p$ is the unitary
\[
 \mathcal U_F:\mathcal R_F\longrightarrow
 L^2([0,\pi]^F,\nu_F),\qquad
 \bigotimes_pe_{n_p}\longmapsto\prod_pq_{n_p}(\theta_p).
\]
It intertwines each local Hecke operation, hence every element of
$\bigotimes_{p\in F}\mathbb C[T_p]$, with multiplication by its
Satake polynomial. The tensor root vector maps to $1$, so its vector
state becomes integration against $\nu_F$. The algebraic integral
lattices are the finite tensor products of the sphere-sum lattices above.
Adjoining a prime uses $v\mapsto v\otimes e_0$ on the tree side and
$f\mapsto f\otimes1$ on the spectral side. These isometric maps
commute with $\mathcal U_F$ and preserve the distinguished states.
They give an operation-compatible finite-stage system.

To specify the occupation-phase comparison, let
$P_p=\mathbb C[a_k,b_k\mid k\ge1]$ be the polynomial oscillator space.
For finitely supported nonnegative integers $r_k,s_k$, give its monomial
$\prod_k a_k^{r_k}b_k^{s_k}$ squared norm
$\prod_k k^{r_k+s_k}r_k!s_k!$, with distinct monomials orthogonal.
Complete this positive inner product to a Hilbert space $\mathcal F_p$.
The charge of this monomial is $m=\sum_k(r_k-s_k)$.
Multiplication by $e^{im\theta}$ on each charge space defines a unitary
$V_\theta$ on $\mathcal F_p$. This is the one-prime occupation action in
Definition~\ref{v4-arith:def:averaging-naive}. The strong integral
$A_p=\int_0^\pi V_\theta\,d\nu_p(\theta)$ is bounded with norm at most one.
It acts on a charge-$m$ monomial by
$c_m=\int_0^\pi e^{im\theta}\,d\nu_p(\theta)$.
For $m=1$, symmetry about $\pi/2$ gives
\[
 c_1=i\int_0^\pi\sin\theta\,d\nu_p(\theta)\in i\mathbb R_{>0}.
\]
Thus this averaging operator has a nonreal eigenvalue. It cannot be
unitarily intertwined with the self-adjoint adjacency operator or
with multiplication by its real spectral coordinate.
This obstruction concerns those specified operators. It does not
exclude a different arithmetic comparison.

The radial transform supplies a complex Hilbert Hecke representation,
an integral cyclic module, and a compatible finite-prime state system.
A comparison with the proposed primitive arithmetic trace class
$\mathcal V^{\mathrm{prim}}$ requires constructing
its arithmetic trace carrier and a map preserving the specified Hecke
operations and trace pairing. Tame residue and Gersten operations
require additional maps on common coefficients.
No full Weil functional or global spectral realization is identified
by the positive root state.


\section{The Chiral Framework over \(\overline{\mathrm{Spec}\,\mathbb{Z}}\)}
\label{v4-arith:kb:beilinson-framework}

The matrices \(S_p\) and their scalar trace recursions define the local
linear constructions above. Global arithmetic realization requires local
representations, distinguished vectors, a restricted tensor-product
construction and an intertwiner with the proposed trace carrier.
Beilinson--Drinfeld factorization on a smooth characteristic-zero curve
does not by itself supply those data over
\(\overline{\mathrm{Spec}\,\mathbb Z}\).

\begin{conjecture}[unramified arithmetic trace realization]
\label{v4-arith:thm:beilinson-honest-reading}
There is an arithmetic dg category $\mathscr C_{\mathrm{arith}}$ over
$\mathbb C$, defined independently of its proposed Hecke realization.
Let $C_\bullet=CC_\bullet(\mathscr C_{\mathrm{arith}})$ be its Hochschild
chain complex, with homological differential of degree $-1$.
The construction specifies a degree $d$, a subspace
$M_{\mathrm{arith}}\subset H_d(C_\bullet)$ containing the primitive
trace class $v_{\mathrm{prim}}=\mathcal V^{\mathrm{prim}}$, and a
normalized linear trace functional
$\tau_{\mathrm{arith}}:M_{\mathrm{arith}}\to\mathbb C$ with
$\tau_{\mathrm{arith}}(v_{\mathrm{prim}})=1$.
It constructs commuting local spherical Hecke actions on this subspace,
with the central element $D_p$ acting as the identity.

For each prime $p$, it also specifies a smooth representation
$\Pi_p$ of $GL_2(\mathbb Q_p)$ with trivial central character,
a Whittaker realization, and the spherical Hecke module
$M_p=\Pi_p^{GL_2(\mathbb Z_p)}$. Its distinguished vector $u_p$
and linear functional $\ell_p:M_p\to\mathbb C$ satisfy
$\ell_p(u_p)=1$. The arithmetic or automorphic construction must
determine these representations, vectors, and functionals.
Set
\[
 M_{\mathrm{aut}}=\bigotimes_p'(M_p,u_p),\qquad
 u=\bigotimes_pu_p,\qquad
 \ell=\bigotimes_p\ell_p.
\]
These are algebraic restricted tensor products over $\mathbb C$.
The normalization makes $\ell$ compatible with adjoining $u_p$.
Each local Hecke algebra acts on its own tensor factor, so the
algebraic restricted tensor product of the local Hecke algebras acts
on $M_{\mathrm{aut}}$.

There is a complex-linear isomorphism
\[
 \Phi:M_{\mathrm{arith}}\longrightarrow M_{\mathrm{aut}},\qquad
 \Phi(v_{\mathrm{prim}})=u,\qquad
 \Phi(hm)=h\Phi(m),\qquad
 \tau_{\mathrm{arith}}(m)=\ell(\Phi(m)).
\]
The identities hold for every $m\in M_{\mathrm{arith}}$ and every
finite-support spherical Hecke element $h$.
In particular, the distinguished Hecke functionals agree:
$\tau_{\mathrm{arith}}(h v_{\mathrm{prim}})=\ell(hu)$.
The oscillator cover additionally admits a complex-linear map
$Q:\bigotimes_p'\mathcal H^{(p)}\to M_{\mathrm{arith}}$ sending its
vacuum to $v_{\mathrm{prim}}$, with complex-linear local maps
$Q_p:\mathcal H^{(p)}\to M_p$ such that
$Q_p(|0\rangle_p)=u_p$ and $\Phi Q=\bigotimes_pQ_p$.
\end{conjecture}

The class $v_{\mathrm{prim}}$ is an element of the arithmetic trace
carrier, whereas $M_{\mathrm{aut}}$ is a representation space.
The conjecture asks for their comparison through $\Phi$, together
with the oscillator maps $Q,Q_p$. It does not identify the class
with that space. All these arithmetic constructions remain open.
The linear functionals above are not assumed positive. A positive
state, Hilbert completion, spectral decomposition, ramified extension,
and archimedean or full Weil comparison each require their own
construction and compatibility theorem.


\section{The Corrected Local Model and Its Consequences}
\label{v4-arith:kb:summary}

The global realization in
Conjecture~\ref{v4-arith:thm:beilinson-honest-reading} remains conjectural.
The unconditional local statements are the divisor composition law,
the distinct Adams monoid action, the symmetric-power trace
recursion for a specified Satake matrix, and the radial spherical
transform of Theorem~\ref{v4-arith:radial:transform}. The latter
intertwines the local Hecke operations and preserves the root state. An arithmetic Hecke action
requires a comparison preserving its carrier and determinant operation.
Neither the oscillator rank nor an equality of scalar traces constructs
that comparison.
